Biological nervous systems are not engineered as circuits; they develop.
A research program that takes development seriously therefore faces a
choice of specification: prescribe the final organization (a designed
network), or prescribe only the materials and the laws and ask whether
organization arises. The ANIMA project takes the second path. Its
charter fixes only the sensory boundary, the output boundary, the
environment and curriculum, finite resource limits, and fundamental
local plasticity laws, and requires that the organism --- a population
of spiking neurons --- develop useful, stable, adaptive internal
organization through interaction, without any prescription of the final
organization \cite{overview,decisionlog}.

Three commitments follow from this framing and structure everything
reported here. First, \emph{minimal specification}: no task-specific
semantic variables, no labels, no reward or credit assignment, no
hidden supervision, and no hand-designed internal architecture. The
organism's ``concepts'' are whatever its plasticity laws and its
curriculum make of its inputs. Second, \emph{locality}: all mechanisms
operate on local state --- membrane potentials, arriving currents,
synaptic weights, and per-neuron budgets. Third, \emph{determinism and
discipline}: every experiment is frozen in writing before
implementation, runs under a single seeded random number generator so
that same-seed reproduction is byte-identical, is gated by a
flag-off identity check against the committed baseline, and preserves
its negative results \cite{overview}.

The research question of Phase I is deliberately modest in form and
demanding in content:

\begin{quote}
\emph{In a minimally specified spiking substrate with local plasticity
and finite resources, what internal organization emerges, what of it
persists, and what are the architectural limits of persistent,
retrievable learning under sequential experience?}
\end{quote}

The answer, as measured rather than assumed, is summarized by the
progression in Figure~\ref{fig:architecture} and developed through the
rest of the paper. The phrase ``architectural limits'' is used
advisedly: Phase I's contribution is not a system that works, but a
map of the causal chain through which each candidate substrate failed,
with the failure point identified at each step. The final step ---
recruiting postsynaptic activity for a genuinely novel pattern that
arrives late in a blocked curriculum --- is a measurement of where the
current architecture stops, not a claim that the class of architectures
is exhausted.

\subsection{Relation to related work}

The substrate is an orthodox leaky-integrate-and-fire (LIF) population
with pairwise additive spike-timing-dependent plasticity (STDP)
\cite{gerstner2002,bipoo1998}. The aspiration --- that competitive
Hebbian learning organizes selective assemblies --- has a long
theoretical lineage \cite{hebb1949,song2000}. ANIMA differs from that
lineage in its discipline rather than its ingredients: organization is
never scored by a task the experimenter cares about; the curriculum is
a fixed, deterministic schedule of synthetic pattern presentations; and
every mechanism claim is tested under flag-off byte-identity control.
The value of the record is therefore mostly negative: it shows, under
these controls, exactly which attractive properties of the ingredients
fail to materialize, and where successive local amendments move the
failure.